%! Linear Combinations of Vectors — Unit 1, Lesson 1.1 — Lesson Plan
\documentclass[10pt]{article}
\usepackage{linalg-article}
\usepackage{linalg-boxes}
\usepackage{graphicx}   % -article does NOT load graphicx; needed for thumbnails/screenshots
\graphicspath{{images/}}
\newcommand{\TallMath}[1]{$\displaystyle #1\rule[-1.4em]{0pt}{3.2em}$}
\newcommand{\vv}{\mathbf{v}}
\newcommand{\ww}{\mathbf{w}}
\newcommand{\UnitNumberName}{Unit 1: Vectors and Matrices \quad}
\newcommand{\LessonNumberName}{Lesson 1.1: Linear Combinations of Vectors}

\begin{document}
\begin{center}
    {\Large\bfseries \CourseName: \SchoolYear \\
    {\small\bfseries \UnitNumberName \LessonNumberName}}
\end{center}

\begin{tcolorbox}[colback=blush, colframe=burgundy, boxrule=0.9pt, arc=2mm,
    left=2mm, right=2mm, top=1.5mm, bottom=1.5mm]
    \textbf{Primary Objective:} Students can build a linear combination $c\,\vv + d\,\ww$ by
    scaling and then adding, and can describe geometrically \emph{which points a set of
    vectors can reach} --- the multiples of one vector fill a line through the origin, and the
    combinations of two non-parallel vectors fill the whole plane --- justifying why parallel
    vectors are the exception.
\end{tcolorbox}

\renewcommand{\arraystretch}{1.4}
\begin{skillbox}[Priority Ideas \& Skills]{goldbox}
\begin{tabularx}{\linewidth}{@{}X|X@{}}
\textbf{Priority Ideas \& Skills} & \textbf{Key Understandings (the ``why'')} \\[2pt]
\hline
\vspace{-6pt}
\begin{itemize}[leftmargin=1.1em, nosep, after=\vspace{2pt}]
  \item Compute a linear combination $c\,\vv+d\,\ww$.
  \item Recognize $\vv+\ww$, $\vv-\ww$, and $c\,\vv$ as special combinations.
  \item Describe the multiples of one vector (a line) and combinations of two (a plane).
  \item Write a target vector as a combination of two given vectors.
\end{itemize}
&
\vspace{-6pt}
\begin{itemize}[leftmargin=1.1em, nosep, after=\vspace{2pt}]
  \item A linear combination does the unit's two operations \emph{at once}: scale, then add.
  \item The reachable points --- the \emph{span} --- have a shape: a line, or the whole plane.
  \item Two non-parallel vectors reach everywhere; parallel vectors stay stuck on one line.
  \item A point's coordinates \emph{are} the weights on $\mathbf{i}$ and $\mathbf{j}$.
\end{itemize}
\end{tabularx}
\end{skillbox}

\begin{skillbox}[{Vocabulary, Concepts, \& Theorems}]{greenbox}
\renewcommand{\arraystretch}{1.3}
\begin{tabularx}{\linewidth}{@{}l X@{}}
\textbf{Linear combination} & \TallMath{c\,\vv + d\,\ww}; scale each vector, then add the results. \\
\textbf{Coefficients (weights)} & The scalars $c,d$ that say how much of each vector to use. \\
\textbf{Span} & All the points you can reach as combinations of the given vectors. \\
\textbf{Parallel vectors} & Two vectors on the same line through the origin --- one is a scalar multiple of the other. \\
\textbf{Standard basis} & \TallMath{\mathbf{i}=\begin{bsmallmatrix} 1 \\ 0 \end{bsmallmatrix},\ \mathbf{j}=\begin{bsmallmatrix} 0 \\ 1 \end{bsmallmatrix}}; every $\begin{bsmallmatrix} x \\ y \end{bsmallmatrix}=x\,\mathbf{i}+y\,\mathbf{j}$. \\
\end{tabularx}
\end{skillbox}

\begin{fixedskillbox}[Activate Prior Knowledge \& Spiral Review]{blush}
    The Warm-Up rehearses the Lesson~1.0 skills this lesson stacks on:
    \textbf{(1)} scaling a vector by a positive and a negative scalar,
    \textbf{(2)} adding two vectors component by component, and
    \textbf{(3)} deciding whether one vector is a \emph{multiple} of another. Debrief item~3
    out loud --- ``is $\begin{bsmallmatrix} 4\\2 \end{bsmallmatrix}$ a multiple of
    $\begin{bsmallmatrix} 2\\1 \end{bsmallmatrix}$?'' --- because being a multiple is exactly
    what puts two vectors on the same line, the idea the whole lesson turns on.
\end{fixedskillbox}

\begin{skillbox}[Hook]{blush}
    Project a delivery robot at the origin with only \textbf{two fixed moves} it can repeat or
    scale: $\vv=\begin{bsmallmatrix} 2\\1 \end{bsmallmatrix}$ and
    $\ww=\begin{bsmallmatrix} 1\\3 \end{bsmallmatrix}$. Ask: \textbf{``Using some amount of
    each move --- $c$ copies of $\vv$ and $d$ copies of $\ww$ --- can the robot land exactly on
    $(5,5)$? On \emph{any} target you name?''} Let students try a few weights before revealing
    $2\vv+\ww=(5,5)$. This frames the lesson's real question: not just \emph{how} to combine
    vectors, but \emph{which points combinations can reach}.
\end{skillbox}

\begin{skillbox}[Lesson]{blush}
\begin{multicols}{2}
\textbf{1. What a linear combination is.} Define $c\,\vv+d\,\ww$ as ``scale each, then add,''
naming $c,d$ the \textbf{coefficients} (weights). Compute
$2\vv+\ww=\begin{bsmallmatrix} 4\\2 \end{bsmallmatrix}+\begin{bsmallmatrix} 1\\3 \end{bsmallmatrix}=\begin{bsmallmatrix} 5\\5 \end{bsmallmatrix}$
with the Hook vectors. Point out the \textbf{special cases}: $\vv+\ww$ ($c=d=1$), $\vv-\ww$
($d=-1$), a lone multiple $c\,\vv$ ($d=0$), and the zero vector ($c=d=0$).

\textbf{2. Multiples of one vector = a line.} Sweep $c$ through all values for a single
$\vv$: every $c\,\vv$ lands on \textbf{one line through the origin}. Use the pre-drawn figure
in the notes; have students read off $2\vv$ and $-\vv$ as points on that line. This is the
\textbf{span} of a single vector.

\columnbreak

\textbf{3. Two vectors = the whole plane (usually).} With two \emph{non-parallel} vectors,
$c\,\vv+d\,\ww$ can reach \textbf{every} point in the plane --- dialing $c$ and $d$ steers you
anywhere. Show $\mathbf{i}=\begin{bsmallmatrix} 1\\0 \end{bsmallmatrix}$,
$\mathbf{j}=\begin{bsmallmatrix} 0\\1 \end{bsmallmatrix}$: since
$\begin{bsmallmatrix} x\\y \end{bsmallmatrix}=x\,\mathbf{i}+y\,\mathbf{j}$, a point's
coordinates \emph{are} its weights. \textbf{The exception:} if $\vv$ and $\ww$ are parallel
(one a multiple of the other), their combinations only refill the same line.

\textbf{4. Finding the weights.} Reverse the question: given a target, find $c,d$. Do
$\begin{bsmallmatrix} 5\\5 \end{bsmallmatrix}=c\vv+d\ww$ by matching components
($2c+d=5,\ c+3d=5$) to recover $c=2,d=1$, tying the answer back to the Hook.
\end{multicols}
\end{skillbox}

\begin{skillbox}[Active Monitoring]{redbox}
Circulate during the notes practice and Tier work. Watch for and correct:
\begin{itemize}[leftmargin=1.2em, nosep]
  \item scaling only one component of a vector, or adding before scaling;
  \item calling two vectors ``different'' when one is actually a multiple of the other (parallel);
  \item claiming two vectors reach the whole plane \emph{without} checking they are non-parallel;
  \item mixing up the weights $c,d$ with the components of the answer vector.
\end{itemize}
Cold-call prompts: ``What does $d=0$ give you?'' ``Point to $2\vv+\ww$ on the grid.''
``Why can't $\begin{bsmallmatrix} 1\\2 \end{bsmallmatrix}$ and $\begin{bsmallmatrix} 2\\4 \end{bsmallmatrix}$
reach $(1,0)$?''
\end{skillbox}

\begin{skillbox}[Group Work \& Differentiation]{redbox}
\textbf{Robot Moves} activity (tiered; mirrors the notes).
\begin{multicols}{3}
\textbf{Tier R --- Remediate.} Given two moves, compute named combinations ($\vv+\ww$, $2\vv$,
$\vv-\ww$) and plot-read where each lands. Rebuilds the ``scale then add'' loop.

\columnbreak
\textbf{Tier A --- Approaching.} Find the weights $c,d$ that hit a target; then test a
\emph{parallel} pair and explain why one target is unreachable.

\columnbreak
\textbf{Tier E --- Extension.} Decide for several vector pairs whether their span is a line or
the plane, and justify the rule (non-parallel $\Rightarrow$ plane) in words.
\end{multicols}
\end{skillbox}

\begin{skillbox}[Individual Work \& Assessment]{redbox}
\textbf{Exit Ticket} (independent, no notes): compute a linear combination; find the weights
$c,d$ for a small target; and a \textbf{conceptual/justification check} --- explain why
$\begin{bsmallmatrix} 2\\4 \end{bsmallmatrix}$ and $\begin{bsmallmatrix} 1\\2 \end{bsmallmatrix}$
cannot reach every point in the plane (they are parallel; span is only a line). Collect and
sort into ``got it / computation slip / span misconception'' to decide whether 1.2 opens with
a quick re-teach.
\end{skillbox}

\begin{skillbox}[Reinforcement \& Extension]{goldbox}
\textbf{Homework:} mixed practice computing combinations, finding weights for a target, and
classifying each vector pair's span as a line or the plane. \textbf{Extension:} a
\emph{data-record} blend --- mixing two base smoothies in amounts $c,d$ shows a linear
combination as a recipe, previewing Unit~1's data theme. \textbf{Preview:} Lesson~1.2,
\emph{Lengths and Angles from Dot Products} --- a single number, the dot product, that measures
how much two vectors point the same way (and recovers length as a special case).
\end{skillbox}
% ── Teacher notes (migrated from the component keys) ──
% One per component, titled for it. They live here rather than in the _key
% files so that every component is the same length blank and keyed.

\begin{teachernote}[Warm-Up]
Item 3 is the pivot for today: ``same multiple of both components'' means the vectors are
\emph{parallel} (same line through the origin). That is exactly the case where two vectors
can \emph{not} reach the whole plane.
\end{teachernote}

\begin{teachernote}[Group Activity]
Tier A item 2 is the key ``interpret'' moment: parallel moves collapse the reachable set to a
line. Tier E item 2 asks for the algebraic reason ($c\vv+dk\vv=(c+dk)\vv$); accept an
equivalent geometric argument (both moves lie on one line, so any combination does too).
\end{teachernote}

\begin{teachernote}[Exit Ticket]
Item 3 checks the central misconception --- assuming \emph{any} two vectors span the plane.
Full credit needs both ``they are parallel / one is a multiple of the other'' \emph{and} the
consequence ``they span only a line.''
\end{teachernote}

\end{document}
